%% annex/p06/infinitary-combinatorics.tex --- proofs for P6 combinatorics

\begin{annexitem}{annex:p6.comb.1}{thm:p06-delta-system}{Delta-system lemma}
\begin{proof}
Let \(\mathcal A\) be an uncountable family of finite sets.  Since
\[
  \mathcal A=\bigcup_{n<\omega}
  \{a\in\mathcal A:\abs a=n\}
\]
and a countable union of countable sets is countable in \(\mathsf{ZFC}\),
there are an \(n<\omega\) and an uncountable
\(\mathcal A_0\subseteq\mathcal A\) consisting of \(n\)-element sets.

We prove the fixed-size assertion by induction on \(n\).  The case \(n=0\)
is vacuous.  Suppose it holds for \(n\), and let \(\mathcal A_0\) be an
uncountable family of \((n+1)\)-element sets.  If \(\mathcal A_0\) has an
uncountable pairwise-disjoint subfamily, that subfamily is a
delta system with root \(\varnothing\).

Otherwise, use Zorn's lemma to choose a maximal pairwise-disjoint
\(\mathcal E\subseteq\mathcal A_0\).  The case assumption makes
\(\mathcal E\) countable, so \(Y=\bigcup\mathcal E\) is countable.
Maximality ensures that every \(a\in\mathcal A_0\) meets \(Y\).  Hence
\[
  \mathcal A_0=
  \bigcup_{x\in Y}\{a\in\mathcal A_0:x\in a\}.
\]
Some \(x\in Y\) therefore belongs to uncountably many members; call that
subfamily \(\mathcal A_x\).  The map
\(a\mapsto a\setminus\{x\}\) is injective on \(\mathcal A_x\), and its
range is an uncountable family of \(n\)-element sets.  By induction it has an
uncountable delta-system subfamily with root \(R_\ast\).  Restoring \(x\)
gives an uncountable delta-system subfamily of \(\mathcal A_x\) with root
\(R_\ast\cup\{x\}\).  This completes the induction.
\end{proof}
\end{annexitem}

\begin{annexitem}{annex:p6.comb.2}{thm:p06-add-ccc}{The Cohen poset is ccc}
\begin{proof}
Let \(A\) be an uncountable family of finite conditions.  The
set \(\{\dom(p):p\in A\}\) is uncountable: a fixed finite domain \(d\)
supports only \(2^{\abs d}\) binary conditions.  Apply the delta-system
lemma to these domains, and choose one condition over each domain in the
resulting uncountable subfamily.  This gives an uncountable
\(A'\subseteq A\) whose domains have a common root \(d_\ast\).  Only finitely
many binary functions on \(d_\ast\) exist, so an uncountable
\(A''\subseteq A'\) has one fixed restriction to \(d_\ast\).  Any two members
of \(A''\) agree on their overlap and their union is
a common strengthening.  Thus no uncountable \(A\) is an antichain.
\end{proof}
\end{annexitem}
